\chapter{Continuous descending fields and mass-decreasing flows}\label{sec:pulltight-deformation}
We give a continuous deformation underlying \cite{MNIndex}, Section 2.8, and the outline in \cite{LiAP}, Section 4.2. Fix a smooth closed connected $M^{n+1}$, $n\geq1$, a mass bound $C>0$, and $K_C,S_C$ from Definition \ref{def:pulltight-mass-cap}. Smooth norms are fixed using a finite atlas or a background connection. The construction below uses a uniformly summable selection of fields, so a locally finite partition of unity is unnecessary.

\begin{lemma}[A negative first variation persists on a neighborhood]\label{lem:local-descending-field}
For every $V\in K_C\setminus S_C$ there are a smooth vector field $X_V$, a number $b_V>0$, and an open neighborhood
\[
 U_V=\{W\in K_C:\delta W(X_V)<-b_V\}
\]
of $V$, disjoint from $S_C$.
\end{lemma}
\begin{proof}
Nonstationarity gives a smooth field with nonzero first variation. Change its sign if necessary to make the variation negative, and choose $b_V$ strictly between zero and $-\delta V(X_V)$. Continuity from Lemma \ref{lem:pulltight-compact-cap} makes $U_V$ open. A stationary varifold has zero first variation on this field, so cannot belong to $U_V$.
\end{proof}

\begin{lemma}[A continuous selection of smooth descending fields]\label{lem:global-descending-field}
There is a map $V\mapsto Y_V$ from $K_C$ to smooth vector fields, continuous in every $C^k$ norm and uniformly bounded in each such norm, with $Y_V=0$ for $V\in S_C$ and $\delta V(Y_V)<0$ for $V\notin S_C$.
\end{lemma}
\begin{proof}
The open subspace $K_C\setminus S_C$ is second countable, so the cover in Lemma \ref{lem:local-descending-field} has a countable subcover $U_i$, with associated fields $X_i$. If this subspace is empty use $Y=0$. Otherwise, repeating sets if the cover is finite, define
\[
 \chi_i(V)=\frac{d(V,K_C\setminus U_i)}{1+d(V,K_C\setminus U_i)},\qquad
 B_i=\max\{1,\|X_i\|_{C^0},\ldots,\|X_i\|_{C^i}\},\qquad
 Y_V=\sum_{i=1}^\infty 2^{-i}\chi_i(V)X_i/B_i.
\]
The complements are nonempty because they contain zero. Each $\chi_i$ is continuous, vanishes off $U_i$, and is positive on $U_i$. For any fixed $k$, the tail starting at $i\geq k$ is bounded in $C^k$ by the summable series $\sum 2^{-i}$. Thus the series converges uniformly in every $C^k$ norm and defines a spatially smooth field; uniform convergence of the continuous finite sums proves continuity of the selection and uniform bounds. On $S_C$ every coefficient vanishes. For other $V$, all terms in $\delta V(Y_V)$ are nonpositive and at least one is strictly negative. Passing first variation through the sum is justified by its $C^1$ convergence and the bound in Lemma \ref{lem:pulltight-compact-cap}.
\end{proof}

\begin{lemma}[A continuous stationarity defect]\label{lem:continuous-stationarity-defect}
For the selection in Lemma \ref{lem:global-descending-field}, set $a(V)=-\delta V(Y_V)$. Then $a$ is continuous and nonnegative on $K_C$, and its zero set is exactly $S_C$. On any nonempty set $\{V:d(V,S_C)\geq\eta\}$ with $\eta>0$, it has a strictly positive minimum.
\end{lemma}
\begin{proof}
Joint continuity of first variation and $C^1$ continuity of $Y$ give continuity of $a$. Its sign and zero set follow from the selection. The set a positive distance from $S_C$ is compact and contains no zero of $a$, so the extreme value theorem gives a positive minimum. This defect is auxiliary and depends on the chosen fields; it is not being identified with a norm of generalized mean curvature.
\end{proof}

\begin{lemma}[Uniform Taylor control along the selected flows]\label{lem:pulltight-flow-taylor}
Let $\varphi_t^V$ be the flow of the fixed field $Y_V$ and let $A_V(t)=\| (\varphi_t^V)_\#V\|(M)$. There is $B\geq1$, independent of $V\in K_C$, such that for $0\leq t\leq1$,
\[
 A'_V(0)=-a(V),\qquad |A''_V(t)|\leq B,\qquad
 A'_V(t)\leq-a(V)+Bt.
\]
\end{lemma}
\begin{proof}
Smooth ODE flow existence and dependence on a vector field in smooth norms are the differential-equation inputs. The fields are complete since $M$ is compact, and have uniform bounds in every smooth norm. Consequently their flows and the needed spatial and time derivatives are uniformly bounded for $0\leq t\leq1$. The rectifiable plane Jacobian, defined for every $P\in G_n(M)$ regardless of rectifiability of $V$, gives
\[
 A_V(t)=\int_{G_n(M)}J_P\varphi_t^V(x)\,dV(x,P).
\]
The first derivative at zero is the first-variation formula. The second time derivative of the Jacobian is bounded uniformly on $K_C\times[0,1]\times G_n(M)$: it is a continuous expression in the flow derivatives, and the flow differential stays invertible with a uniform inverse bound on this interval. Integrating this bound against mass at most $C$ gives $|A''_V|\leq B$, after enlarging $B$ to at least one. Integration in time yields the last inequality.
\end{proof}

\begin{proposition}[A continuous time choice with a positive mass drop]\label{prop:pulltight-time-choice}
Fix $0<\tau_0\leq1$ and put
\[
 \tau(V)=\min\{\tau_0,a(V)/(2B)\},\qquad
 D(V)=\tfrac12 a(V)\tau(V),\qquad
 \mathscr H(s,V)=(\varphi_{s\tau(V)}^V)_\#V.
\]
Then $\mathscr H:[0,1]\times K_C\to K_C$ is continuous, fixes every stationary varifold, starts at the identity, and has nonincreasing mass in $s$. At its endpoint,
\[
 \|\mathscr H(1,V)\|(M)\leq\|V\|(M)-D(V).
\]
The continuous function $D$ vanishes exactly on $S_C$.
\end{proposition}
\begin{proof}
For $0\leq t\leq\tau(V)$, Lemma \ref{lem:pulltight-flow-taylor} gives $A'_V(t)\leq-a(V)/2$. Integrate from zero to $s\tau(V)$ to obtain monotonicity and the endpoint estimate. For stationary $V$, both the field and the chosen time are zero. For nonstationary $V$, both $a$ and $\tau$ are positive, so $D>0$. Continuity of the fields and times gives $C^1$ continuity of the diffeomorphisms in $(s,V)$. Lemma \ref{lem:pulltight-pushforward-continuity} gives continuity of $\mathscr H$. Its mass bound shows it stays in $K_C$.
\end{proof}

\begin{corollary}[Uniform loss away from stationary varifolds]\label{cor:pulltight-uniform-loss}
For every $\eta>0$ such that $K_\eta=\{V\in K_C:d(V,S_C)\geq\eta\}$ is nonempty, there is $d_\eta>0$ with
\[
 \|\mathscr H(1,V)\|(M)\leq\|V\|(M)-d_\eta\quad(V\in K_\eta).
\]
In particular, if the mass losses of a sequence tend to zero, every convergent subsequence of the original varifolds has stationary limit.
\end{corollary}
\begin{proof}
Take $d_\eta=\min_{K_\eta}D>0$ using compactness and positivity from Proposition \ref{prop:pulltight-time-choice}. If $V_j\to V$ and the actual losses tend to zero, the bound $0\leq D(V_j)\leq\|V_j\|(M)-\|\mathscr H(1,V_j)\|(M)$ implies $D(V)=0$. Thus $V\in S_C$. A fixed deformation and its fixed positive defect are used here; arbitrarily small losses for a changing sequence of unrelated deformations would not yield this conclusion.
\end{proof}

\begin{lemma}[Uniform closeness to the identity]\label{lem:pulltight-near-identity}
The choice of $\tau_0$ in Proposition \ref{prop:pulltight-time-choice} can make the deformation arbitrarily close to the identity in $\mathbf F_V$, uniformly on $[0,1]\times K_C$. For every integral mod-two cycle $T$ with mass at most $C$, the corresponding diffeomorphisms also satisfy
\[
 \mathcal F((\varphi_t^{|T|})_\#T-T)\leq A t\mathbf M(T),\qquad 0\leq t\leq1,
\]
with $A$ independent of $T$.
\end{lemma}
\begin{proof}
Uniform smooth bounds give $d(\varphi_t(x),x)\leq A_1t$, displacement of the tangent plane at most $A_1t$, and $|J_P\varphi_t-1|\leq A_1t$, with a fixed smooth metric on $G_n(M)$. For a test function of absolute value and Lipschitz constant at most one, subtract the pushforward and original integrals to bound their difference by $A_2Ct$. Taking the supremum gives the varifold estimate. For cycles, use the external mod-two homotopy formula for $h(u,x)=\varphi_u^{|T|}(x)$ on $[0,t]\times M$. Its time derivative is uniformly bounded and its spatial derivatives are bounded, hence the swept $(n+1)$-chain has mass at most $At\mathbf M(T)$. Its boundary is the required difference because $\partial T=0$. Taking $t\leq\tau_0$ makes the sum of the two bounds arbitrarily small. No mass-distance bound follows from this argument; Example \ref{ex:translated-slab-mass-distance} rules one out in general.
\end{proof}

\begin{theorem}[A continuous pull-tight homotopy on boundary cycles]\label{thm:continuous-pulltight-homotopy}
Let $M^{n+1}$ be smooth, closed and connected, $n\geq1$, and let $C,\varepsilon>0$. On the mass-at-most-$C$ part of $\mathcal Z_n^0(M;\mathbb Z_2)$ there is a jointly $\mathbf F$-continuous homotopy $H(s,T)$ which starts at $T$, fixes $T$ whenever $|T|$ is stationary, has nonincreasing mass in $s$, and satisfies
\[
 \mathbf M(H(1,T))\leq\mathbf M(T)-D(|T|),\qquad
 \mathbf F(H(s,T),T)\leq\varepsilon,
\]
where $D$ is continuous on $K_C$ and positive precisely off $S_C$. This proves the continuous deformation stated in \cite{MNIndex}, Section 2.8, from the current operations and ODE inputs specified above.
\end{theorem}
\begin{proof}
Choose $\tau_0$ using Lemma \ref{lem:pulltight-near-identity} and define $H(s,T)=(\varphi_{s\tau(|T|)}^{|T|})_\#T$. Diffeomorphic pushforward preserves integral mod-two cycles and commutes with boundary. In particular if $T=\partial U$, its image is the boundary of the image of $U$, so remains in the boundary component. The Jacobian formula with no cancellation identifies $|H(s,T)|=\mathscr H(s,|T|)$. Proposition \ref{prop:pulltight-time-choice} supplies mass monotonicity, the drop, and the fixed-point properties. Lemma \ref{lem:pulltight-near-identity} supplies the distance estimate. Finally $\mathbf F$ convergence of cycles gives weak convergence of their varifolds, hence $C^1$ convergence of the selected diffeomorphisms. Lemma \ref{lem:pulltight-pushforward-continuity} proves joint $\mathbf F$ continuity. No discrete mass-fineness assertion is used or proved.
\end{proof}
